\chapter{Client Pain Points}

\section{Recurring Operational Pain Points}

Across sectors, common pain points:
\begin{itemize}
    \item \textbf{Manual internal transport}: parts, pallets, medications, room deliveries, mail -- slow, repetitive, fatigue-prone. In hospitals, staff spend much time on transport \cite{researchgate_hospital,mir_healthcare}.
    \item \textbf{Repetitive/hazardous inspection}: manual gauge reading, thermal/leak checks -- slow, infrequent, dangerous.
    \item \textbf{Labor scarcity and peaks}: seasonal swings, turnover hard to manage.
    \item \textbf{Data/coverage gaps}: inconsistent inspections, fixed-camera blind spots.
\end{itemize}

\section{Overarching Constraint: Cost-to-Labor Ratio}

These pains are shared with high-wage markets, but the economics differ. Tunisian wages are much lower, so a robot for pure labor substitution has long payback. International sheets ignore this.

Thus, focus on sectors where ROI compares to something other than wages:
\begin{itemize}
    \item \textbf{Inspection}: avoided incident cost (accident, downtime) can exceed robot cost.
    \item \textbf{Universities}: compared to imported research platform cost, not salary.
\end{itemize}

\section{Financial and Procurement Pain}

Local buyers are capex-averse; procurement is slow (public sector). Imported hardware adds currency risk, customs delays, lead-time friction. Thin-margin operators cannot absorb large capex. These constraints push toward service-based models \cite{builtin_raas}.